%%%%%%%%%%%%%%%%%%%%%%%%%%%%%%%%%%%%%%%%%%%%%%%%%%%%%%%%%%%%
% 6.11 DESIGN THEORY
% The Composition of Advanced Mathematics
%%%%%%%%%%%%%%%%%%%%%%%%%%%%%%%%%%%%%%%%%%%%%%%%%%%%%%%%%%%%

\section{Design Theory}
\label{sec:design-theory}

\prerequisite{
Basic counting, combinations, finite sets, graph theory, incidence
structures, finite fields, elementary number theory, and basic linear
algebra.
}

\learningobjectives{
By the end of this section, the reader should be able to:
\begin{itemize}
    \item explain the purpose of combinatorial design theory;
    \item define incidence structures, points, blocks, and incidence;
    \item define \(t\)-designs and their parameters;
    \item work with balanced incomplete block designs;
    \item derive the fundamental parameter equations for BIBDs;
    \item determine necessary divisibility conditions for designs;
    \item define and recognize Steiner systems;
    \item work with Steiner triple systems;
    \item understand the existence condition for Steiner triple systems;
    \item define symmetric designs;
    \item understand projective planes as combinatorial designs;
    \item construct finite projective planes from finite fields;
    \item define affine planes;
    \item understand resolvable designs and parallel classes;
    \item recognize Latin squares and mutually orthogonal Latin squares;
    \item connect design theory with finite geometry, coding theory,
          statistics, graph theory, and algebraic combinatorics;
    \item use incidence matrices to study designs;
    \item recognize difference sets and group-developed designs;
    \item understand pairwise balanced designs and covering designs;
    \item distinguish exact designs from packing and covering problems.
\end{itemize}
}

%%%%%%%%%%%%%%%%%%%%%%%%%%%%%%%%%%%%%%%%%%%%%%%%%%%%%%%%%%%%
\subsection{What Is Design Theory?}
\label{subsec:what-is-design-theory}
%%%%%%%%%%%%%%%%%%%%%%%%%%%%%%%%%%%%%%%%%%%%%%%%%%%%%%%%%%%%

Design theory studies finite collections of subsets arranged so that
elements occur together with prescribed regularity.

The central question is:

\[
\boxed{
\text{Can a finite set be organized into highly balanced families of
subsets?}
}
\]

A combinatorial design consists of:

\begin{itemize}
    \item a set of points;
    \item a collection of blocks;
    \item regular incidence relationships between points and blocks.
\end{itemize}

Design theory is simultaneously combinatorial, algebraic, geometric,
and statistical.

%%%%%%%%%%%%%%%%%%%%%%%%%%%%%%%%%%%%%%%%%%%%%%%%%%%%%%%%%%%%
\subsection{Incidence Structures}
\label{subsec:incidence-structures-design}
%%%%%%%%%%%%%%%%%%%%%%%%%%%%%%%%%%%%%%%%%%%%%%%%%%%%%%%%%%%%

\begin{definitionbox}{Incidence Structure}
An incidence structure is a pair

\[
\boxed{
(X,\mathcal{B}),
}
\]

where \(X\) is a finite set of points and

\[
\mathcal{B}
\]

is a collection of subsets of \(X\), called blocks.
\end{definitionbox}

A point \(x\) is incident with a block \(B\) when

\[
\boxed{
x\in B.
}
\]

Thus ordinary set membership becomes the incidence relation.

%%%%%%%%%%%%%%%%%%%%%%%%%%%%%%%%%%%%%%%%%%%%%%%%%%%%%%%%%%%%
\subsection{Basic Parameters}
\label{subsec:design-basic-parameters}
%%%%%%%%%%%%%%%%%%%%%%%%%%%%%%%%%%%%%%%%%%%%%%%%%%%%%%%%%%%%

Let

\[
|X|=v.
\]

The symbol

\[
v
\]

denotes the number of points.

Let

\[
|\mathcal{B}|=b.
\]

The symbol

\[
b
\]

denotes the number of blocks.

If every block has size \(k\), then

\[
\boxed{
|B|=k
}
\]

for every

\[
B\in\mathcal{B}.
\]

If every point occurs in exactly \(r\) blocks, then \(r\) is called the
replication number.

%%%%%%%%%%%%%%%%%%%%%%%%%%%%%%%%%%%%%%%%%%%%%%%%%%%%%%%%%%%%
\subsection{\(t\)-Designs}
\label{subsec:t-designs}
%%%%%%%%%%%%%%%%%%%%%%%%%%%%%%%%%%%%%%%%%%%%%%%%%%%%%%%%%%%%

\begin{definitionbox}{\(t\)-Design}
A \(t\)-\((v,k,\lambda)\) design is an incidence structure

\[
(X,\mathcal{B})
\]

such that:

\begin{itemize}
    \item \(|X|=v\);
    \item every block has size \(k\);
    \item every \(t\)-element subset of \(X\) occurs in exactly
          \(\lambda\) blocks.
\end{itemize}
\end{definitionbox}

The Greek letter

\[
\lambda
\]

is called ``lambda.''

The notation is read

\[
\boxed{
t\text{-}(v,k,\lambda)\text{ design}.
}
\]

%%%%%%%%%%%%%%%%%%%%%%%%%%%%%%%%%%%%%%%%%%%%%%%%%%%%%%%%%%%%
\subsection{Interpreting the Parameters}
\label{subsec:interpreting-design-parameters}
%%%%%%%%%%%%%%%%%%%%%%%%%%%%%%%%%%%%%%%%%%%%%%%%%%%%%%%%%%%%

In a

\[
t\text{-}(v,k,\lambda)
\]

design:

\[
v
=
\text{number of points},
\]

\[
k
=
\text{points per block},
\]

\[
t
=
\text{size of subsets whose incidence is controlled},
\]

and

\[
\lambda
=
\text{number of blocks containing each }t\text{-subset}.
\]

Increasing \(t\) imposes stronger regularity.

%%%%%%%%%%%%%%%%%%%%%%%%%%%%%%%%%%%%%%%%%%%%%%%%%%%%%%%%%%%%
\subsection{Balanced Incomplete Block Designs}
\label{subsec:bibd}
%%%%%%%%%%%%%%%%%%%%%%%%%%%%%%%%%%%%%%%%%%%%%%%%%%%%%%%%%%%%

The most classical designs are \(2\)-designs.

\begin{definitionbox}{Balanced Incomplete Block Design}
A balanced incomplete block design, abbreviated BIBD, is a
\(2\)-\((v,k,\lambda)\) design.
\end{definitionbox}

Thus every pair of distinct points occurs together in exactly

\[
\lambda
\]

blocks.

The word ``incomplete'' reflects that

\[
k<v,
\]

so blocks do not contain every point.

%%%%%%%%%%%%%%%%%%%%%%%%%%%%%%%%%%%%%%%%%%%%%%%%%%%%%%%%%%%%
\subsection{BIBD Parameters}
\label{subsec:bibd-parameters}
%%%%%%%%%%%%%%%%%%%%%%%%%%%%%%%%%%%%%%%%%%%%%%%%%%%%%%%%%%%%

A BIBD is often described by five parameters:

\[
\boxed{
(v,b,r,k,\lambda).
}
\]

These mean:

\[
v
=
\text{number of points},
\]

\[
b
=
\text{number of blocks},
\]

\[
r
=
\text{blocks containing each point},
\]

\[
k
=
\text{points in each block},
\]

\[
\lambda
=
\text{blocks containing each pair of points}.
\]

These parameters cannot be chosen independently.

%%%%%%%%%%%%%%%%%%%%%%%%%%%%%%%%%%%%%%%%%%%%%%%%%%%%%%%%%%%%
\subsection{Counting Point--Block Incidences}
\label{subsec:counting-point-block-incidences}
%%%%%%%%%%%%%%%%%%%%%%%%%%%%%%%%%%%%%%%%%%%%%%%%%%%%%%%%%%%%

Count incidence pairs

\[
(x,B)
\]

with

\[
x\in B.
\]

Counting by points gives

\[
vr.
\]

Counting by blocks gives

\[
bk.
\]

Therefore,

\[
\boxed{
vr=bk.
}
\]

This is the first fundamental BIBD identity.

%%%%%%%%%%%%%%%%%%%%%%%%%%%%%%%%%%%%%%%%%%%%%%%%%%%%%%%%%%%%
\subsection{Counting Pairs Through a Point}
\label{subsec:counting-pairs-through-point}
%%%%%%%%%%%%%%%%%%%%%%%%%%%%%%%%%%%%%%%%%%%%%%%%%%%%%%%%%%%%

Fix a point \(x\).

The point \(x\) occurs in exactly \(r\) blocks.

Each such block contains

\[
k-1
\]

additional points.

Thus the number of pairs

\[
(y,B)
\]

with

\[
y\neq x
\]

and

\[
x,y\in B
\]

is

\[
r(k-1).
\]

On the other hand, there are

\[
v-1
\]

possible choices of \(y\), and each pair \(\{x,y\}\) occurs in
\(\lambda\) blocks.

Therefore,

\[
\boxed{
r(k-1)
=
\lambda(v-1).
}
\]

This is the second fundamental BIBD identity.

%%%%%%%%%%%%%%%%%%%%%%%%%%%%%%%%%%%%%%%%%%%%%%%%%%%%%%%%%%%%
\subsection{Number of Blocks}
\label{subsec:bibd-number-blocks}
%%%%%%%%%%%%%%%%%%%%%%%%%%%%%%%%%%%%%%%%%%%%%%%%%%%%%%%%%%%%

From

\[
r(k-1)=\lambda(v-1),
\]

we obtain

\[
\boxed{
r
=
\frac{\lambda(v-1)}{k-1}.
}
\]

Using

\[
vr=bk,
\]

we then obtain

\[
\boxed{
b
=
\frac{\lambda v(v-1)}
{k(k-1)}.
}
\]

Thus \(r\) and \(b\) are determined by \(v,k,\lambda\).

%%%%%%%%%%%%%%%%%%%%%%%%%%%%%%%%%%%%%%%%%%%%%%%%%%%%%%%%%%%%
\subsection{Worked Example: BIBD Parameters}
\label{subsec:worked-bibd-parameters}
%%%%%%%%%%%%%%%%%%%%%%%%%%%%%%%%%%%%%%%%%%%%%%%%%%%%%%%%%%%%

\begin{workedexamplebox}{A \(2\)-\((7,3,1)\) Design}

Suppose

\[
v=7,
\qquad
k=3,
\qquad
\lambda=1.
\]

Then

\[
r
=
\frac{1(7-1)}{3-1}
=
3.
\]

Next,

\[
b
=
\frac{7\cdot3}{3}
=
7.
\]

\begin{answerbox}
\[
\boxed{
(v,b,r,k,\lambda)
=
(7,7,3,3,1).
}
\]
\end{answerbox}

\end{workedexamplebox}

%%%%%%%%%%%%%%%%%%%%%%%%%%%%%%%%%%%%%%%%%%%%%%%%%%%%%%%%%%%%
\subsection{Necessary Divisibility Conditions}
\label{subsec:design-divisibility-conditions}
%%%%%%%%%%%%%%%%%%%%%%%%%%%%%%%%%%%%%%%%%%%%%%%%%%%%%%%%%%%%

For a BIBD to exist,

\[
r
=
\frac{\lambda(v-1)}{k-1}
\]

must be an integer.

Therefore,

\[
\boxed{
k-1
\mid
\lambda(v-1).
}
\]

Similarly,

\[
b
=
\frac{\lambda v(v-1)}
{k(k-1)}
\]

must be an integer.

Hence,

\[
\boxed{
k(k-1)
\mid
\lambda v(v-1).
}
\]

These are necessary conditions.

They are not always sufficient.

%%%%%%%%%%%%%%%%%%%%%%%%%%%%%%%%%%%%%%%%%%%%%%%%%%%%%%%%%%%%
\subsection{Fisher's Inequality}
\label{subsec:fishers-inequality}
%%%%%%%%%%%%%%%%%%%%%%%%%%%%%%%%%%%%%%%%%%%%%%%%%%%%%%%%%%%%

\begin{theorem}[Fisher's Inequality]
For every nontrivial BIBD,

\[
\boxed{
b\geq v.
}
\]
\end{theorem}

Thus the number of blocks is at least the number of points.

The theorem can be proved using linear algebra and the incidence matrix
of the design.

%%%%%%%%%%%%%%%%%%%%%%%%%%%%%%%%%%%%%%%%%%%%%%%%%%%%%%%%%%%%
\subsection{Symmetric Designs}
\label{subsec:symmetric-designs}
%%%%%%%%%%%%%%%%%%%%%%%%%%%%%%%%%%%%%%%%%%%%%%%%%%%%%%%%%%%%

\begin{definitionbox}{Symmetric Design}
A BIBD is symmetric if

\[
\boxed{
b=v.
}
\]
\end{definitionbox}

Using

\[
vr=bk,
\]

a symmetric design also satisfies

\[
\boxed{
r=k.
}
\]

Symmetric designs possess strong duality between points and blocks.

%%%%%%%%%%%%%%%%%%%%%%%%%%%%%%%%%%%%%%%%%%%%%%%%%%%%%%%%%%%%
\subsection{Block Intersections in Symmetric Designs}
\label{subsec:symmetric-block-intersections}
%%%%%%%%%%%%%%%%%%%%%%%%%%%%%%%%%%%%%%%%%%%%%%%%%%%%%%%%%%%%

A remarkable property of a symmetric

\[
2\text{-}(v,k,\lambda)
\]

design is that every two distinct blocks intersect in exactly

\[
\boxed{
\lambda
}
\]

points.

Thus the pairwise regularity among points is mirrored by pairwise
regularity among blocks.

%%%%%%%%%%%%%%%%%%%%%%%%%%%%%%%%%%%%%%%%%%%%%%%%%%%%%%%%%%%%
\subsection{Dual Designs}
\label{subsec:dual-designs}
%%%%%%%%%%%%%%%%%%%%%%%%%%%%%%%%%%%%%%%%%%%%%%%%%%%%%%%%%%%%

Given an incidence structure, one may interchange the roles of points
and blocks.

The original blocks become points.

The original points become blocks.

Incidence is preserved.

This produces the dual incidence structure.

For symmetric designs, the dual has the same basic parameters as the
original design.

%%%%%%%%%%%%%%%%%%%%%%%%%%%%%%%%%%%%%%%%%%%%%%%%%%%%%%%%%%%%
\subsection{Steiner Systems}
\label{subsec:steiner-systems}
%%%%%%%%%%%%%%%%%%%%%%%%%%%%%%%%%%%%%%%%%%%%%%%%%%%%%%%%%%%%

\begin{definitionbox}{Steiner System}
A Steiner system

\[
\boxed{
S(t,k,v)
}
\]

is a \(t\)-\((v,k,1)\) design.
\end{definitionbox}

Thus every \(t\)-element subset occurs in exactly one block.

Steiner systems represent perfect combinatorial decompositions.

%%%%%%%%%%%%%%%%%%%%%%%%%%%%%%%%%%%%%%%%%%%%%%%%%%%%%%%%%%%%
\subsection{Steiner Triple Systems}
\label{subsec:steiner-triple-systems}
%%%%%%%%%%%%%%%%%%%%%%%%%%%%%%%%%%%%%%%%%%%%%%%%%%%%%%%%%%%%

A Steiner triple system is a Steiner system

\[
\boxed{
S(2,3,v).
}
\]

Thus:

\begin{itemize}
    \item there are \(v\) points;
    \item each block contains \(3\) points;
    \item every pair of points occurs in exactly one block.
\end{itemize}

It is commonly abbreviated

\[
\boxed{
\operatorname{STS}(v).
}
\]

%%%%%%%%%%%%%%%%%%%%%%%%%%%%%%%%%%%%%%%%%%%%%%%%%%%%%%%%%%%%
\subsection{Parameters of a Steiner Triple System}
\label{subsec:sts-parameters}
%%%%%%%%%%%%%%%%%%%%%%%%%%%%%%%%%%%%%%%%%%%%%%%%%%%%%%%%%%%%

For

\[
k=3,
\qquad
\lambda=1,
\]

we obtain

\[
r
=
\frac{v-1}{2}
\]

and

\[
b
=
\frac{v(v-1)}{6}.
\]

Thus,

\[
\boxed{
r=\frac{v-1}{2},
\qquad
b=\frac{v(v-1)}{6}.
}
\]

These quantities must be integers.

%%%%%%%%%%%%%%%%%%%%%%%%%%%%%%%%%%%%%%%%%%%%%%%%%%%%%%%%%%%%
\subsection{Existence of Steiner Triple Systems}
\label{subsec:sts-existence}
%%%%%%%%%%%%%%%%%%%%%%%%%%%%%%%%%%%%%%%%%%%%%%%%%%%%%%%%%%%%

A Steiner triple system exists exactly for certain values of \(v\).

\begin{theorem}
A Steiner triple system of order \(v\) exists if and only if

\[
\boxed{
v\equiv1\text{ or }3\pmod6.
}
\]
\end{theorem}

The notation

\[
v\equiv1\pmod6
\]

is read ``\(v\) is congruent to \(1\) modulo \(6\).''

%%%%%%%%%%%%%%%%%%%%%%%%%%%%%%%%%%%%%%%%%%%%%%%%%%%%%%%%%%%%
\subsection{Why the Congruence Condition Appears}
\label{subsec:sts-congruence}
%%%%%%%%%%%%%%%%%%%%%%%%%%%%%%%%%%%%%%%%%%%%%%%%%%%%%%%%%%%%

We require

\[
r=\frac{v-1}{2}
\]

to be an integer.

Therefore \(v\) must be odd.

We also require

\[
b=\frac{v(v-1)}{6}
\]

to be an integer.

Combining the divisibility requirements yields

\[
\boxed{
v\equiv1
\quad\text{or}\quad
3
\pmod6.
}
\]

The nontrivial part of the theorem is showing that these necessary
conditions are also sufficient.

%%%%%%%%%%%%%%%%%%%%%%%%%%%%%%%%%%%%%%%%%%%%%%%%%%%%%%%%%%%%
\subsection{The Fano Plane}
\label{subsec:fano-plane-design}
%%%%%%%%%%%%%%%%%%%%%%%%%%%%%%%%%%%%%%%%%%%%%%%%%%%%%%%%%%%%

The smallest nontrivial Steiner triple system is

\[
\operatorname{STS}(7).
\]

It is also the finite projective plane of order \(2\).

Its parameters are

\[
\boxed{
2\text{-}(7,3,1).
}
\]

It contains:

\[
7
\]

points and

\[
7
\]

blocks.

Every block contains \(3\) points.

Every pair of points lies in exactly one block.

%%%%%%%%%%%%%%%%%%%%%%%%%%%%%%%%%%%%%%%%%%%%%%%%%%%%%%%%%%%%
\subsection{Worked Example: A Fano Plane Block System}
\label{subsec:worked-fano-blocks}
%%%%%%%%%%%%%%%%%%%%%%%%%%%%%%%%%%%%%%%%%%%%%%%%%%%%%%%%%%%%

\begin{workedexamplebox}{Seven Blocks on Seven Points}

Let

\[
X=\{1,2,3,4,5,6,7\}.
\]

Consider the blocks

\[
\{1,2,3\},
\]

\[
\{1,4,5\},
\]

\[
\{1,6,7\},
\]

\[
\{2,4,6\},
\]

\[
\{2,5,7\},
\]

\[
\{3,4,7\},
\]

\[
\{3,5,6\}.
\]

Every pair of points occurs in exactly one block.

\begin{answerbox}
\[
\boxed{
(X,\mathcal{B})
\text{ is a }2\text{-}(7,3,1)\text{ design}.
}
\]
\end{answerbox}

\end{workedexamplebox}

%%%%%%%%%%%%%%%%%%%%%%%%%%%%%%%%%%%%%%%%%%%%%%%%%%%%%%%%%%%%
\subsection{Projective Planes}
\label{subsec:projective-planes-design}
%%%%%%%%%%%%%%%%%%%%%%%%%%%%%%%%%%%%%%%%%%%%%%%%%%%%%%%%%%%%

A finite projective plane is an incidence structure satisfying:

\begin{itemize}
    \item every pair of distinct points lies on exactly one line;
    \item every pair of distinct lines meets in exactly one point;
    \item there exist four points with no three collinear.
\end{itemize}

The lines serve as blocks.

A projective plane of order \(q\) has highly regular parameters.

%%%%%%%%%%%%%%%%%%%%%%%%%%%%%%%%%%%%%%%%%%%%%%%%%%%%%%%%%%%%
\subsection{Parameters of a Projective Plane}
\label{subsec:projective-plane-parameters}
%%%%%%%%%%%%%%%%%%%%%%%%%%%%%%%%%%%%%%%%%%%%%%%%%%%%%%%%%%%%

A projective plane of order \(q\) has

\[
\boxed{
v=q^2+q+1
}
\]

points and the same number of lines.

Thus,

\[
\boxed{
b=q^2+q+1.
}
\]

Every line contains

\[
\boxed{
q+1
}
\]

points, and every point lies on

\[
q+1
\]

lines.

Therefore a projective plane of order \(q\) is a symmetric design with

\[
\boxed{
2\text{-}(q^2+q+1,q+1,1).
}
\]

%%%%%%%%%%%%%%%%%%%%%%%%%%%%%%%%%%%%%%%%%%%%%%%%%%%%%%%%%%%%
\subsection{Finite Projective Planes from Finite Fields}
\label{subsec:projective-planes-finite-fields}
%%%%%%%%%%%%%%%%%%%%%%%%%%%%%%%%%%%%%%%%%%%%%%%%%%%%%%%%%%%%

When \(q\) is a prime power, a projective plane of order \(q\) can be
constructed from the vector space

\[
\mathbb{F}_q^3.
\]

Points correspond to one-dimensional subspaces.

Lines correspond to two-dimensional subspaces.

Incidence is given by subspace containment.

This construction is denoted

\[
\boxed{
\operatorname{PG}(2,q).
}
\]

The notation is read ``projective geometry of dimension \(2\) over
\(\mathbb{F}_q\).''

%%%%%%%%%%%%%%%%%%%%%%%%%%%%%%%%%%%%%%%%%%%%%%%%%%%%%%%%%%%%
\subsection{Counting Projective Points}
\label{subsec:counting-projective-points}
%%%%%%%%%%%%%%%%%%%%%%%%%%%%%%%%%%%%%%%%%%%%%%%%%%%%%%%%%%%%

There are

\[
q^3-1
\]

nonzero vectors in

\[
\mathbb{F}_q^3.
\]

Each one-dimensional subspace contains

\[
q-1
\]

nonzero vectors.

Therefore the number of one-dimensional subspaces is

\[
\frac{q^3-1}{q-1}.
\]

Factor:

\[
q^3-1
=
(q-1)(q^2+q+1).
\]

Hence,

\[
\boxed{
v=q^2+q+1.
}
\]

%%%%%%%%%%%%%%%%%%%%%%%%%%%%%%%%%%%%%%%%%%%%%%%%%%%%%%%%%%%%
\subsection{Affine Planes}
\label{subsec:affine-planes}
%%%%%%%%%%%%%%%%%%%%%%%%%%%%%%%%%%%%%%%%%%%%%%%%%%%%%%%%%%%%

An affine plane of order \(q\) has

\[
\boxed{
q^2
}
\]

points.

Each line contains

\[
q
\]

points.

Through each point pass

\[
q+1
\]

lines.

Lines fall into parallel classes.

An affine plane of order \(q\) forms a

\[
\boxed{
2\text{-}(q^2,q,1)
}
\]

design.

%%%%%%%%%%%%%%%%%%%%%%%%%%%%%%%%%%%%%%%%%%%%%%%%%%%%%%%%%%%%
\subsection{Parallel Classes}
\label{subsec:parallel-classes}
%%%%%%%%%%%%%%%%%%%%%%%%%%%%%%%%%%%%%%%%%%%%%%%%%%%%%%%%%%%%

A parallel class is a collection of pairwise disjoint blocks whose
union is the entire point set.

Thus every point appears in exactly one block of the class.

If every block has size \(k\), then a parallel class contains

\[
\boxed{
\frac{v}{k}
}
\]

blocks.

%%%%%%%%%%%%%%%%%%%%%%%%%%%%%%%%%%%%%%%%%%%%%%%%%%%%%%%%%%%%
\subsection{Resolvable Designs}
\label{subsec:resolvable-designs}
%%%%%%%%%%%%%%%%%%%%%%%%%%%%%%%%%%%%%%%%%%%%%%%%%%%%%%%%%%%%

\begin{definitionbox}{Resolvable Design}
A design is resolvable if its block set can be partitioned into parallel
classes.
\end{definitionbox}

Resolvable designs are useful in scheduling because each parallel class
can represent one round in which every participant appears exactly
once.

%%%%%%%%%%%%%%%%%%%%%%%%%%%%%%%%%%%%%%%%%%%%%%%%%%%%%%%%%%%%
\subsection{Kirkman's Schoolgirl Problem}
\label{subsec:kirkman-schoolgirl}
%%%%%%%%%%%%%%%%%%%%%%%%%%%%%%%%%%%%%%%%%%%%%%%%%%%%%%%%%%%%

A famous design-theory problem asks whether \(15\) schoolgirls can walk
in groups of \(3\) for \(7\) days so that every pair walks together
exactly once.

The problem asks for a resolvable Steiner triple system of order \(15\).

Each day is a parallel class containing

\[
5
\]

triples.

The complete schedule uses

\[
7
\]

parallel classes.

%%%%%%%%%%%%%%%%%%%%%%%%%%%%%%%%%%%%%%%%%%%%%%%%%%%%%%%%%%%%
\subsection{Pairwise Balanced Designs}
\label{subsec:pairwise-balanced-designs}
%%%%%%%%%%%%%%%%%%%%%%%%%%%%%%%%%%%%%%%%%%%%%%%%%%%%%%%%%%%%

A pairwise balanced design relaxes the condition that every block have
the same size.

\begin{definitionbox}{Pairwise Balanced Design}
A pairwise balanced design is an incidence structure in which every pair
of distinct points occurs in exactly one block, while block sizes may
vary.
\end{definitionbox}

If the allowed block sizes belong to a set \(K\), the design is often
written

\[
\boxed{
\operatorname{PBD}(v,K).
}
\]

%%%%%%%%%%%%%%%%%%%%%%%%%%%%%%%%%%%%%%%%%%%%%%%%%%%%%%%%%%%%
\subsection{Packing Designs}
\label{subsec:packing-designs}
%%%%%%%%%%%%%%%%%%%%%%%%%%%%%%%%%%%%%%%%%%%%%%%%%%%%%%%%%%%%

A packing design relaxes the exact incidence requirement.

A \(t\)-\((v,k,\lambda)\) packing requires every \(t\)-subset to occur in
\emph{at most} \(\lambda\) blocks.

The goal is often to maximize the number of blocks.

Thus packing theory is an extremal version of design theory.

%%%%%%%%%%%%%%%%%%%%%%%%%%%%%%%%%%%%%%%%%%%%%%%%%%%%%%%%%%%%
\subsection{Covering Designs}
\label{subsec:covering-designs}
%%%%%%%%%%%%%%%%%%%%%%%%%%%%%%%%%%%%%%%%%%%%%%%%%%%%%%%%%%%%

A covering design requires every \(t\)-subset to occur in
\emph{at least} \(\lambda\) blocks.

The goal is often to minimize the number of blocks.

Thus:

\[
\boxed{
\text{packing}
=
\text{avoid excessive repetition},
}
\]

while

\[
\boxed{
\text{covering}
=
\text{guarantee sufficient coverage}.
}
\]

%%%%%%%%%%%%%%%%%%%%%%%%%%%%%%%%%%%%%%%%%%%%%%%%%%%%%%%%%%%%
\subsection{Covering Numbers}
\label{subsec:covering-numbers}
%%%%%%%%%%%%%%%%%%%%%%%%%%%%%%%%%%%%%%%%%%%%%%%%%%%%%%%%%%%%

The covering number

\[
\boxed{
C(v,k,t)
}
\]

is the minimum number of \(k\)-element blocks required so that every
\(t\)-element subset of a \(v\)-element set is contained in at least one
block.

Exact covering numbers are difficult to determine in many cases.

%%%%%%%%%%%%%%%%%%%%%%%%%%%%%%%%%%%%%%%%%%%%%%%%%%%%%%%%%%%%
\subsection{Latin Squares}
\label{subsec:latin-squares}
%%%%%%%%%%%%%%%%%%%%%%%%%%%%%%%%%%%%%%%%%%%%%%%%%%%%%%%%%%%%

\begin{definitionbox}{Latin Square}
A Latin square of order \(n\) is an \(n\times n\) array filled with
\(n\) symbols such that every symbol occurs exactly once in every row
and exactly once in every column.
\end{definitionbox}

For example, a cyclic Latin square of order \(3\) is

\[
\begin{array}{c|ccc}
 &1&2&3\\
\hline
1&1&2&3\\
2&2&3&1\\
3&3&1&2
\end{array}
\]

Latin squares are closely related to experimental design and finite
geometry.

%%%%%%%%%%%%%%%%%%%%%%%%%%%%%%%%%%%%%%%%%%%%%%%%%%%%%%%%%%%%
\subsection{Latin Squares from Groups}
\label{subsec:latin-squares-groups}
%%%%%%%%%%%%%%%%%%%%%%%%%%%%%%%%%%%%%%%%%%%%%%%%%%%%%%%%%%%%

The Cayley table of a finite group is a Latin square.

Fixing a row corresponds to left multiplication by a group element,
which is a bijection.

Similarly, each column contains every group element exactly once.

Thus algebra naturally produces Latin squares.

%%%%%%%%%%%%%%%%%%%%%%%%%%%%%%%%%%%%%%%%%%%%%%%%%%%%%%%%%%%%
\subsection{Orthogonal Latin Squares}
\label{subsec:orthogonal-latin-squares}
%%%%%%%%%%%%%%%%%%%%%%%%%%%%%%%%%%%%%%%%%%%%%%%%%%%%%%%%%%%%

Two Latin squares of order \(n\) are orthogonal if superimposing them
produces every ordered pair of symbols exactly once.

Such a pair contains

\[
n^2
\]

cells and

\[
n^2
\]

possible ordered symbol pairs.

%%%%%%%%%%%%%%%%%%%%%%%%%%%%%%%%%%%%%%%%%%%%%%%%%%%%%%%%%%%%
\subsection{Mutually Orthogonal Latin Squares}
\label{subsec:mols}
%%%%%%%%%%%%%%%%%%%%%%%%%%%%%%%%%%%%%%%%%%%%%%%%%%%%%%%%%%%%

A family of Latin squares is mutually orthogonal if every distinct pair
is orthogonal.

The abbreviation

\[
\boxed{
\text{MOLS}
}
\]

stands for mutually orthogonal Latin squares.

For a prime power \(q\), there exist

\[
\boxed{
q-1
}
\]

mutually orthogonal Latin squares of order \(q\).

Such a family is called complete.

%%%%%%%%%%%%%%%%%%%%%%%%%%%%%%%%%%%%%%%%%%%%%%%%%%%%%%%%%%%%
\subsection{MOLS and Finite Geometry}
\label{subsec:mols-finite-geometry}
%%%%%%%%%%%%%%%%%%%%%%%%%%%%%%%%%%%%%%%%%%%%%%%%%%%%%%%%%%%%

A complete set of

\[
q-1
\]

MOLS of order \(q\) is closely related to an affine plane of order
\(q\).

Thus the following structures are deeply connected:

\[
\boxed{
\text{finite fields}
\longleftrightarrow
\text{affine planes}
\longleftrightarrow
\text{MOLS}
\longleftrightarrow
\text{block designs}.
}
\]

%%%%%%%%%%%%%%%%%%%%%%%%%%%%%%%%%%%%%%%%%%%%%%%%%%%%%%%%%%%%
\subsection{Designs in Experimental Statistics}
\label{subsec:designs-experimental-statistics}
%%%%%%%%%%%%%%%%%%%%%%%%%%%%%%%%%%%%%%%%%%%%%%%%%%%%%%%%%%%%

The term ``block design'' arose partly from statistical experimental
design.

Suppose \(v\) treatments are to be tested in blocks of size \(k\).

A BIBD arranges treatments so that every pair occurs together equally
often.

This balances pairwise comparisons and reduces confounding caused by
block effects.

%%%%%%%%%%%%%%%%%%%%%%%%%%%%%%%%%%%%%%%%%%%%%%%%%%%%%%%%%%%%
\subsection{Replication and Balance}
\label{subsec:replication-balance}
%%%%%%%%%%%%%%%%%%%%%%%%%%%%%%%%%%%%%%%%%%%%%%%%%%%%%%%%%%%%

In a BIBD:

\[
r
\]

controls replication of individual treatments.

The parameter

\[
\lambda
\]

controls pairwise concurrence.

The identities

\[
vr=bk
\]

and

\[
r(k-1)=\lambda(v-1)
\]

ensure global consistency between these forms of balance.

%%%%%%%%%%%%%%%%%%%%%%%%%%%%%%%%%%%%%%%%%%%%%%%%%%%%%%%%%%%%
\subsection{Incidence Matrices}
\label{subsec:design-incidence-matrices}
%%%%%%%%%%%%%%%%%%%%%%%%%%%%%%%%%%%%%%%%%%%%%%%%%%%%%%%%%%%%

Let a design have points

\[
x_1,\ldots,x_v
\]

and blocks

\[
B_1,\ldots,B_b.
\]

Its incidence matrix is the \(v\times b\) matrix

\[
N=(n_{ij}),
\]

where

\[
\boxed{
n_{ij}
=
\begin{cases}
1,&x_i\in B_j,\\
0,&x_i\notin B_j.
\end{cases}
}
\]

The matrix converts incidence information into linear algebra.

%%%%%%%%%%%%%%%%%%%%%%%%%%%%%%%%%%%%%%%%%%%%%%%%%%%%%%%%%%%%
\subsection{Matrix Identity for BIBDs}
\label{subsec:bibd-matrix-identity}
%%%%%%%%%%%%%%%%%%%%%%%%%%%%%%%%%%%%%%%%%%%%%%%%%%%%%%%%%%%%

For a BIBD, the matrix

\[
NN^T
\]

has:

\begin{itemize}
    \item diagonal entries \(r\);
    \item off-diagonal entries \(\lambda\).
\end{itemize}

Therefore,

\[
\boxed{
NN^T
=
(r-\lambda)I
+
\lambda J,
}
\]

where:

\[
I
\]

is the identity matrix, and

\[
J
\]

is the all-ones matrix.

This identity is fundamental in the linear-algebraic study of designs.

%%%%%%%%%%%%%%%%%%%%%%%%%%%%%%%%%%%%%%%%%%%%%%%%%%%%%%%%%%%%
\subsection{Fisher's Inequality from Linear Algebra}
\label{subsec:fisher-linear-algebra}
%%%%%%%%%%%%%%%%%%%%%%%%%%%%%%%%%%%%%%%%%%%%%%%%%%%%%%%%%%%%

For a nontrivial BIBD,

\[
r>\lambda.
\]

Therefore,

\[
(r-\lambda)I+\lambda J
\]

is nonsingular.

Hence

\[
NN^T
\]

has rank \(v\).

Since

\[
\operatorname{rank}(NN^T)
\leq
\operatorname{rank}(N)
\leq
b,
\]

we obtain

\[
v\leq b.
\]

Thus,

\[
\boxed{
b\geq v.
}
\]

This proves Fisher's inequality.

%%%%%%%%%%%%%%%%%%%%%%%%%%%%%%%%%%%%%%%%%%%%%%%%%%%%%%%%%%%%
\subsection{Block Intersection Graphs}
\label{subsec:block-intersection-graphs}
%%%%%%%%%%%%%%%%%%%%%%%%%%%%%%%%%%%%%%%%%%%%%%%%%%%%%%%%%%%%

A design can be converted into a graph.

Create one vertex for each block.

Join two vertices when the corresponding blocks intersect.

This produces a block intersection graph.

Graph invariants can then reveal structural information about the
design.

%%%%%%%%%%%%%%%%%%%%%%%%%%%%%%%%%%%%%%%%%%%%%%%%%%%%%%%%%%%%
\subsection{Incidence Graphs}
\label{subsec:design-incidence-graphs}
%%%%%%%%%%%%%%%%%%%%%%%%%%%%%%%%%%%%%%%%%%%%%%%%%%%%%%%%%%%%

Every incidence structure determines a bipartite graph.

One part consists of points.

The other part consists of blocks.

Connect point \(x\) to block \(B\) whenever

\[
x\in B.
\]

The resulting graph is called the incidence graph or Levi graph.

%%%%%%%%%%%%%%%%%%%%%%%%%%%%%%%%%%%%%%%%%%%%%%%%%%%%%%%%%%%%
\subsection{The Fano Plane Incidence Graph}
\label{subsec:fano-incidence-graph}
%%%%%%%%%%%%%%%%%%%%%%%%%%%%%%%%%%%%%%%%%%%%%%%%%%%%%%%%%%%%

The incidence graph of the Fano plane has:

\[
7
\]

point vertices and

\[
7
\]

block vertices.

Thus it has

\[
14
\]

vertices.

Every vertex has degree \(3\).

This graph is known as the Heawood graph.

It illustrates the strong connection between design theory and regular
bipartite graphs.

%%%%%%%%%%%%%%%%%%%%%%%%%%%%%%%%%%%%%%%%%%%%%%%%%%%%%%%%%%%%
\subsection{Automorphisms of Designs}
\label{subsec:design-automorphisms}
%%%%%%%%%%%%%%%%%%%%%%%%%%%%%%%%%%%%%%%%%%%%%%%%%%%%%%%%%%%%

An automorphism of a design is a permutation of the point set that maps
blocks to blocks.

The automorphisms form a group.

Thus group actions can be used to study:

\begin{itemize}
    \item symmetry;
    \item classification;
    \item equivalent designs;
    \item orbit structure;
    \item efficient constructions.
\end{itemize}

Highly symmetric designs often arise from finite groups or finite
geometries.

%%%%%%%%%%%%%%%%%%%%%%%%%%%%%%%%%%%%%%%%%%%%%%%%%%%%%%%%%%%%
\subsection{Difference Sets}
\label{subsec:difference-sets}
%%%%%%%%%%%%%%%%%%%%%%%%%%%%%%%%%%%%%%%%%%%%%%%%%%%%%%%%%%%%

Let \(G\) be a finite group of order \(v\), written additively when
convenient.

A subset

\[
D\subseteq G
\]

with

\[
|D|=k
\]

is a difference set if every nonidentity element of \(G\) occurs exactly

\[
\lambda
\]

times among differences

\[
d_1-d_2
\]

with

\[
d_1,d_2\in D,
\qquad
d_1\neq d_2.
\]

Such a set is called a

\[
\boxed{
(v,k,\lambda)\text{ difference set}.
}
\]

%%%%%%%%%%%%%%%%%%%%%%%%%%%%%%%%%%%%%%%%%%%%%%%%%%%%%%%%%%%%
\subsection{Designs from Difference Sets}
\label{subsec:designs-from-difference-sets}
%%%%%%%%%%%%%%%%%%%%%%%%%%%%%%%%%%%%%%%%%%%%%%%%%%%%%%%%%%%%

Translate a difference set \(D\) by every group element:

\[
D+g
=
\{d+g:d\in D\}.
\]

The collection of all translates produces the blocks of a symmetric
design.

Thus difference sets provide an algebraic construction method for
designs.

%%%%%%%%%%%%%%%%%%%%%%%%%%%%%%%%%%%%%%%%%%%%%%%%%%%%%%%%%%%%
\subsection{Difference Set Parameter Equation}
\label{subsec:difference-set-equation}
%%%%%%%%%%%%%%%%%%%%%%%%%%%%%%%%%%%%%%%%%%%%%%%%%%%%%%%%%%%%

There are

\[
k(k-1)
\]

ordered differences

\[
d_1-d_2
\]

with

\[
d_1\neq d_2.
\]

There are

\[
v-1
\]

nonidentity group elements, each occurring exactly \(\lambda\) times.

Therefore,

\[
\boxed{
k(k-1)
=
\lambda(v-1).
}
\]

This is the same relation that appears in symmetric BIBDs.

%%%%%%%%%%%%%%%%%%%%%%%%%%%%%%%%%%%%%%%%%%%%%%%%%%%%%%%%%%%%
\subsection{Hadamard Designs Preview}
\label{subsec:hadamard-designs-preview}
%%%%%%%%%%%%%%%%%%%%%%%%%%%%%%%%%%%%%%%%%%%%%%%%%%%%%%%%%%%%

Hadamard matrices can be used to construct important symmetric designs.

A Hadamard matrix is a square matrix with entries

\[
\pm1
\]

whose rows are pairwise orthogonal.

If \(H\) has order \(n\), then

\[
\boxed{
HH^T=nI.
}
\]

Hadamard matrices connect:

\begin{itemize}
    \item design theory;
    \item coding theory;
    \item finite geometry;
    \item signal processing.
\end{itemize}

%%%%%%%%%%%%%%%%%%%%%%%%%%%%%%%%%%%%%%%%%%%%%%%%%%%%%%%%%%%%
\subsection{Codes from Designs}
\label{subsec:codes-from-designs}
%%%%%%%%%%%%%%%%%%%%%%%%%%%%%%%%%%%%%%%%%%%%%%%%%%%%%%%%%%%%

The rows or columns of an incidence matrix may be interpreted as binary
vectors.

Working over

\[
\mathbb{F}_2,
\]

their span defines a binary linear code.

Thus a design may generate a code whose algebraic properties reflect
the combinatorial regularity of the design.

%%%%%%%%%%%%%%%%%%%%%%%%%%%%%%%%%%%%%%%%%%%%%%%%%%%%%%%%%%%%
\subsection{Designs from Codes}
\label{subsec:designs-from-codes}
%%%%%%%%%%%%%%%%%%%%%%%%%%%%%%%%%%%%%%%%%%%%%%%%%%%%%%%%%%%%

The relationship also works in the opposite direction.

Supports of codewords of a fixed weight may form blocks of a
combinatorial design.

This creates a deep connection between:

\[
\boxed{
\text{designs}
+
\text{finite fields}
+
\text{error-correcting codes}.
}
\]

%%%%%%%%%%%%%%%%%%%%%%%%%%%%%%%%%%%%%%%%%%%%%%%%%%%%%%%%%%%%
\subsection{Assmus--Mattson Phenomenon Preview}
\label{subsec:assmus-mattson-preview}
%%%%%%%%%%%%%%%%%%%%%%%%%%%%%%%%%%%%%%%%%%%%%%%%%%%%%%%%%%%%

Certain highly structured linear codes have codewords whose supports
form \(t\)-designs.

The Assmus--Mattson theorem gives conditions under which this occurs.

This is an advanced example of algebraic regularity producing
combinatorial regularity.

%%%%%%%%%%%%%%%%%%%%%%%%%%%%%%%%%%%%%%%%%%%%%%%%%%%%%%%%%%%%
\subsection{Orthogonal Arrays}
\label{subsec:orthogonal-arrays}
%%%%%%%%%%%%%%%%%%%%%%%%%%%%%%%%%%%%%%%%%%%%%%%%%%%%%%%%%%%%

An orthogonal array is a rectangular array with controlled frequency of
symbol patterns in selected columns.

A common notation is

\[
\boxed{
\operatorname{OA}(N,k,s,t).
}
\]

Roughly:

\[
N
=
\text{number of rows},
\]

\[
k
=
\text{number of columns},
\]

\[
s
=
\text{number of symbols},
\]

and

\[
t
=
\text{strength}.
\]

Every selection of \(t\) columns contains all \(t\)-tuples equally
often.

%%%%%%%%%%%%%%%%%%%%%%%%%%%%%%%%%%%%%%%%%%%%%%%%%%%%%%%%%%%%
\subsection{Orthogonal Arrays and Experimental Design}
\label{subsec:orthogonal-arrays-experiments}
%%%%%%%%%%%%%%%%%%%%%%%%%%%%%%%%%%%%%%%%%%%%%%%%%%%%%%%%%%%%

Orthogonal arrays provide efficient experimental layouts.

Instead of testing every possible factor combination, one can choose a
balanced subset that preserves controlled interactions.

They are widely used in:

\begin{itemize}
    \item industrial experiments;
    \item quality engineering;
    \item coding theory;
    \item computer experiments.
\end{itemize}

%%%%%%%%%%%%%%%%%%%%%%%%%%%%%%%%%%%%%%%%%%%%%%%%%%%%%%%%%%%%
\subsection{Design Existence Problems}
\label{subsec:design-existence-problems}
%%%%%%%%%%%%%%%%%%%%%%%%%%%%%%%%%%%%%%%%%%%%%%%%%%%%%%%%%%%%

A central problem is determining whether a design with specified
parameters exists.

Necessary conditions are often obtained from:

\begin{itemize}
    \item divisibility;
    \item counting;
    \item linear algebra;
    \item number theory;
    \item eigenvalue restrictions.
\end{itemize}

Sufficiency may require difficult explicit constructions or deep
existence theorems.

%%%%%%%%%%%%%%%%%%%%%%%%%%%%%%%%%%%%%%%%%%%%%%%%%%%%%%%%%%%%
\subsection{Necessary Does Not Mean Sufficient}
\label{subsec:necessary-not-sufficient-design}
%%%%%%%%%%%%%%%%%%%%%%%%%%%%%%%%%%%%%%%%%%%%%%%%%%%%%%%%%%%%

Suppose parameters satisfy

\[
vr=bk
\]

and

\[
r(k-1)=\lambda(v-1).
\]

This does not automatically guarantee that a BIBD exists.

These equations are necessary arithmetic conditions.

A valid block structure must still be constructed.

Thus design existence contains both:

\[
\boxed{
\text{arithmetic feasibility}
}
\]

and

\[
\boxed{
\text{combinatorial realizability}.
}
\]

%%%%%%%%%%%%%%%%%%%%%%%%%%%%%%%%%%%%%%%%%%%%%%%%%%%%%%%%%%%%
\subsection{Wilson-Type Existence Philosophy}
\label{subsec:wilson-existence-preview}
%%%%%%%%%%%%%%%%%%%%%%%%%%%%%%%%%%%%%%%%%%%%%%%%%%%%%%%%%%%%

Modern design theory contains powerful asymptotic existence theorems.

A general phenomenon is that, once the natural divisibility conditions
are satisfied and the parameters are sufficiently large, appropriate
designs often exist.

This idea appears in major existence results for block designs,
pairwise balanced designs, and related decomposition problems.

%%%%%%%%%%%%%%%%%%%%%%%%%%%%%%%%%%%%%%%%%%%%%%%%%%%%%%%%%%%%
\subsection{Graph Decompositions and Designs}
\label{subsec:graph-decompositions-designs}
%%%%%%%%%%%%%%%%%%%%%%%%%%%%%%%%%%%%%%%%%%%%%%%%%%%%%%%%%%%%

A Steiner system

\[
S(2,k,v)
\]

can be interpreted as a decomposition of the edge set of \(K_v\) into
copies of \(K_k\).

Every pair of points is an edge.

Every block of size \(k\) contributes a copy of \(K_k\).

Because every pair occurs once, the copies partition

\[
E(K_v).
\]

Thus,

\[
\boxed{
\text{block designs}
\longleftrightarrow
\text{graph decompositions}.
}
\]

%%%%%%%%%%%%%%%%%%%%%%%%%%%%%%%%%%%%%%%%%%%%%%%%%%%%%%%%%%%%
\subsection{Steiner Triple Systems as Triangle Decompositions}
\label{subsec:sts-triangle-decomposition}
%%%%%%%%%%%%%%%%%%%%%%%%%%%%%%%%%%%%%%%%%%%%%%%%%%%%%%%%%%%%

For an

\[
S(2,3,v),
\]

each block is a triple.

Each triple represents a triangle.

Therefore a Steiner triple system is exactly a decomposition of

\[
K_v
\]

into edge-disjoint triangles.

This gives another explanation for the divisibility conditions:

\[
\binom v2
\]

must be divisible by

\[
3.
\]

%%%%%%%%%%%%%%%%%%%%%%%%%%%%%%%%%%%%%%%%%%%%%%%%%%%%%%%%%%%%
\subsection{Designs and Hypergraphs}
\label{subsec:designs-hypergraphs}
%%%%%%%%%%%%%%%%%%%%%%%%%%%%%%%%%%%%%%%%%%%%%%%%%%%%%%%%%%%%

A \(k\)-uniform block design may be viewed as a \(k\)-uniform
hypergraph.

The point set becomes the vertex set.

The blocks become hyperedges.

A \(t\)-design imposes uniform codegree:

\[
\boxed{
\text{every }t\text{-set lies in exactly }\lambda\text{ hyperedges}.
}
\]

Thus design theory is a highly regular part of hypergraph theory.

%%%%%%%%%%%%%%%%%%%%%%%%%%%%%%%%%%%%%%%%%%%%%%%%%%%%%%%%%%%%
\subsection{Designs and Ramsey Theory}
\label{subsec:designs-ramsey}
%%%%%%%%%%%%%%%%%%%%%%%%%%%%%%%%%%%%%%%%%%%%%%%%%%%%%%%%%%%%

Ramsey theory seeks unavoidable structure.

Design theory seeks deliberately constructed regular structure.

Both investigate how local relationships interact with global
organization.

Some highly symmetric designs also produce strong Ramsey-type
constructions.

%%%%%%%%%%%%%%%%%%%%%%%%%%%%%%%%%%%%%%%%%%%%%%%%%%%%%%%%%%%%
\subsection{Designs and Matroids}
\label{subsec:designs-matroids}
%%%%%%%%%%%%%%%%%%%%%%%%%%%%%%%%%%%%%%%%%%%%%%%%%%%%%%%%%%%%

Finite geometries underlying many designs also produce matroids.

For example, points of

\[
\operatorname{PG}(n,q)
\]

carry dependence relations inherited from vector spaces over
\(\mathbb{F}_q\).

Thus projective designs connect naturally with representable matroids,
which will be developed in the next section.

%%%%%%%%%%%%%%%%%%%%%%%%%%%%%%%%%%%%%%%%%%%%%%%%%%%%%%%%%%%%
\subsection{Common Errors}
\label{subsec:design-theory-common-errors}
%%%%%%%%%%%%%%%%%%%%%%%%%%%%%%%%%%%%%%%%%%%%%%%%%%%%%%%%%%%%

\subsubsection{Confusing Points and Blocks}

Points are the underlying elements.

Blocks are subsets of the point set.

%%%%%%%%%%%%%%%%%%%%%%%%%%%%%%%%%%%%%%%%%%%%%%%%%%%%%%%%%%%%
\subsubsection{Confusing \(t\) and \(k\)}

In a

\[
t\text{-}(v,k,\lambda)
\]

design:

\[
k
\]

is the block size, while

\[
t
\]

is the size of subsets whose incidence is controlled.

%%%%%%%%%%%%%%%%%%%%%%%%%%%%%%%%%%%%%%%%%%%%%%%%%%%%%%%%%%%%
\subsubsection{Assuming Every \(t\)-Subset Is a Block}

A \(t\)-subset need not itself be a block.

It must be contained in exactly

\[
\lambda
\]

blocks.

%%%%%%%%%%%%%%%%%%%%%%%%%%%%%%%%%%%%%%%%%%%%%%%%%%%%%%%%%%%%
\subsubsection{Forgetting Replication}

In a BIBD, every point appears in the same number \(r\) of blocks.

This follows from the design conditions and should not be treated as an
independent arbitrary parameter.

%%%%%%%%%%%%%%%%%%%%%%%%%%%%%%%%%%%%%%%%%%%%%%%%%%%%%%%%%%%%
\subsubsection{Ignoring Divisibility Conditions}

If

\[
r
\]

or

\[
b
\]

is not an integer, the proposed BIBD cannot exist.

%%%%%%%%%%%%%%%%%%%%%%%%%%%%%%%%%%%%%%%%%%%%%%%%%%%%%%%%%%%%
\subsubsection{Assuming Divisibility Guarantees Existence}

Arithmetic conditions are necessary but not always sufficient.

A combinatorial construction is still required.

%%%%%%%%%%%%%%%%%%%%%%%%%%%%%%%%%%%%%%%%%%%%%%%%%%%%%%%%%%%%
\subsubsection{Confusing Steiner Systems and General Designs}

A Steiner system has

\[
\lambda=1.
\]

General \(t\)-designs may have larger \(\lambda\).

%%%%%%%%%%%%%%%%%%%%%%%%%%%%%%%%%%%%%%%%%%%%%%%%%%%%%%%%%%%%
\subsubsection{Confusing Projective and Affine Planes}

A projective plane has no parallel lines.

Any two lines intersect exactly once.

An affine plane contains parallel classes.

%%%%%%%%%%%%%%%%%%%%%%%%%%%%%%%%%%%%%%%%%%%%%%%%%%%%%%%%%%%%
\subsubsection{Confusing Latin Squares with Block Designs}

Latin squares are arrays.

Block designs are collections of subsets.

They are related, but they are not the same object.

%%%%%%%%%%%%%%%%%%%%%%%%%%%%%%%%%%%%%%%%%%%%%%%%%%%%%%%%%%%%
\subsubsection{Confusing Packing and Covering Conditions}

A packing requires

\[
\text{at most }\lambda
\]

occurrences.

A covering requires

\[
\text{at least }\lambda.
\]

A design requires exactly \(\lambda\).

%%%%%%%%%%%%%%%%%%%%%%%%%%%%%%%%%%%%%%%%%%%%%%%%%%%%%%%%%%%%
\subsection{Applications}
\label{subsec:design-theory-applications}
%%%%%%%%%%%%%%%%%%%%%%%%%%%%%%%%%%%%%%%%%%%%%%%%%%%%%%%%%%%%

\begin{applicationbox}

\textbf{Experimental Design.}
Balanced block designs allow treatments to be compared fairly when not
all treatments can be tested simultaneously.

\medskip

\textbf{Coding Theory.}
Incidence matrices and finite geometries generate structured
error-correcting codes.

\medskip

\textbf{Finite Geometry.}
Projective and affine planes are among the most important examples of
combinatorial designs.

\medskip

\textbf{Scheduling.}
Resolvable designs organize participants into balanced groups across
multiple rounds.

\medskip

\textbf{Networking.}
Block designs can construct balanced communication and testing
patterns.

\medskip

\textbf{Cryptography.}
Finite incidence structures and difference sets appear in sequence
design and combinatorial cryptographic constructions.

\medskip

\textbf{Software and Hardware Testing.}
Covering arrays and orthogonal arrays reduce the number of tests while
ensuring systematic interaction coverage.

\medskip

\textbf{Graph Theory.}
Steiner systems correspond to decompositions of complete graphs into
smaller complete subgraphs.

\medskip

\textbf{Algebraic Combinatorics.}
Automorphism groups, association schemes, difference sets, and
incidence matrices reveal algebraic structure inside designs.

\end{applicationbox}

%%%%%%%%%%%%%%%%%%%%%%%%%%%%%%%%%%%%%%%%%%%%%%%%%%%%%%%%%%%%
\subsection{Exercises}
\label{subsec:design-theory-exercises}
%%%%%%%%%%%%%%%%%%%%%%%%%%%%%%%%%%%%%%%%%%%%%%%%%%%%%%%%%%%%

\begin{exercise}[1]
Define an incidence structure.
\end{exercise}

\begin{exercise}[1]
In a block design, what is a point?
\end{exercise}

\begin{exercise}[1]
In a block design, what is a block?
\end{exercise}

\begin{exercise}[1]
Define a

\[
t\text{-}(v,k,\lambda)
\]

design.
\end{exercise}

\begin{exercise}[1]
Define a BIBD.
\end{exercise}

\begin{exercise}[1]
Explain the meanings of

\[
v,b,r,k,\lambda.
\]
\end{exercise}

\begin{exercise}[1]
State the identity

\[
vr=bk.
\]
\end{exercise}

\begin{exercise}[1]
State the identity involving

\[
r(k-1).
\]
\end{exercise}

\begin{exercise}[1]
Define a Steiner system.
\end{exercise}

\begin{exercise}[1]
Define a Steiner triple system.
\end{exercise}

\begin{exercise}[1]
What are the parameters of the Fano plane as a design?
\end{exercise}

\begin{exercise}[1]
Define a symmetric design.
\end{exercise}

\begin{exercise}[2]
Derive

\[
vr=bk
\]

by double counting point-block incidences.
\end{exercise}

\begin{exercise}[2]
Derive

\[
r(k-1)
=
\lambda(v-1).
\]
\end{exercise}

\begin{exercise}[2]
Derive

\[
r
=
\frac{\lambda(v-1)}{k-1}.
\]
\end{exercise}

\begin{exercise}[2]
Derive

\[
b
=
\frac{\lambda v(v-1)}
{k(k-1)}.
\]
\end{exercise}

\begin{exercise}[2]
Determine \(r\) and \(b\) for a hypothetical

\[
2\text{-}(13,4,1)
\]

design.
\end{exercise}

\begin{exercise}[2]
Determine whether the arithmetic parameters permit a

\[
2\text{-}(10,3,1)
\]

design.
\end{exercise}

\begin{exercise}[2]
For a Steiner triple system, derive

\[
r=\frac{v-1}{2}.
\]
\end{exercise}

\begin{exercise}[2]
For a Steiner triple system, derive

\[
b=\frac{v(v-1)}6.
\]
\end{exercise}

\begin{exercise}[2]
Use divisibility conditions to show that an \(\operatorname{STS}(v)\)
can exist only when

\[
v\equiv1
\quad\text{or}\quad
3
\pmod6.
\]
\end{exercise}

\begin{exercise}[2]
Verify directly that the seven blocks listed for the Fano plane contain
every pair exactly once.
\end{exercise}

\begin{exercise}[2]
How many points does a projective plane of order \(3\) contain?
\end{exercise}

\begin{exercise}[2]
How many lines does a projective plane of order \(3\) contain?
\end{exercise}

\begin{exercise}[2]
How many points lie on each line of a projective plane of order \(3\)?
\end{exercise}

\begin{exercise}[2]
Find the design parameters of an affine plane of order \(4\).
\end{exercise}

\begin{exercise}[2]
Define a parallel class.
\end{exercise}

\begin{exercise}[2]
Define a resolvable design.
\end{exercise}

\begin{exercise}[2]
Construct a Latin square of order \(4\) using addition modulo \(4\).
\end{exercise}

\begin{exercise}[2]
Explain why the Cayley table of a finite group is a Latin square.
\end{exercise}

\begin{exercise}[3]
Prove Fisher's inequality using the incidence matrix identity

\[
NN^T
=
(r-\lambda)I+\lambda J.
\]
\end{exercise}

\begin{exercise}[3]
Show that a symmetric BIBD satisfies

\[
r=k.
\]
\end{exercise}

\begin{exercise}[3]
Prove that any two distinct blocks of a symmetric BIBD intersect in
exactly \(\lambda\) points.
\end{exercise}

\begin{exercise}[3]
Show that the Fano plane is self-dual.
\end{exercise}

\begin{exercise}[3]
Interpret a Steiner triple system as a triangle decomposition of
\(K_v\).
\end{exercise}

\begin{exercise}[3]
Using the graph-decomposition interpretation, derive necessary
divisibility conditions for an \(\operatorname{STS}(v)\).
\end{exercise}

\begin{exercise}[3]
Construct the incidence matrix of the Fano plane.
\end{exercise}

\begin{exercise}[3]
Verify

\[
NN^T
=
2I+J
\]

for the Fano plane incidence matrix.
\end{exercise}

\begin{exercise}[3]
Show that the incidence graph of the Fano plane is \(3\)-regular.
\end{exercise}

\begin{exercise}[3]
Count the one-dimensional subspaces of

\[
\mathbb{F}_q^3
\]

and derive

\[
q^2+q+1.
\]
\end{exercise}

\begin{exercise}[3]
Explain why two-dimensional subspaces of

\[
\mathbb{F}_q^3
\]

serve as projective lines.
\end{exercise}

\begin{exercise}[3]
Prove that every projective line in

\[
\operatorname{PG}(2,q)
\]

contains \(q+1\) points.
\end{exercise}

\begin{exercise}[3]
Show that a complete set of \(q-1\) MOLS of order \(q\) produces an
affine plane of order \(q\).
\end{exercise}

\begin{exercise}[3]
Let \(D\) be a \((v,k,\lambda)\) difference set. Prove

\[
k(k-1)=\lambda(v-1).
\]
\end{exercise}

\begin{exercise}[3]
Explain how translating a difference set produces a symmetric design.
\end{exercise}

\begin{exercise}[3]
Compare a \(t\)-design, a \(t\)-packing, and a \(t\)-covering.
\end{exercise}

\begin{exercise}[4]
Study the existence theorem for Steiner triple systems.
\end{exercise}

\begin{exercise}[4]
Investigate Steiner systems

\[
S(3,4,v)
\]

and determine their necessary divisibility conditions.
\end{exercise}

\begin{exercise}[4]
Study finite projective planes and the open problem concerning projective
planes of non-prime-power order.
\end{exercise}

\begin{exercise}[4]
Investigate the Bruck--Ryser--Chowla theorem.
\end{exercise}

\begin{exercise}[4]
Study Hadamard matrices and their relationship with symmetric block
designs.
\end{exercise}

\begin{exercise}[4]
Investigate difference sets in cyclic groups.
\end{exercise}

\begin{exercise}[4]
Study the relationship between orthogonal arrays and error-correcting
codes.
\end{exercise}

\begin{exercise}[4]
Investigate covering arrays used in software testing.
\end{exercise}

\begin{exercise}[4]
Study Wilson's existence theory for pairwise balanced designs.
\end{exercise}

\begin{exercise}[4]
Investigate automorphism groups of the Fano plane and their relation to
finite linear groups.
\end{exercise}

%%%%%%%%%%%%%%%%%%%%%%%%%%%%%%%%%%%%%%%%%%%%%%%%%%%%%%%%%%%%
\subsection{Conceptual Summary}
\label{subsec:design-theory-summary}
%%%%%%%%%%%%%%%%%%%%%%%%%%%%%%%%%%%%%%%%%%%%%%%%%%%%%%%%%%%%

Design theory studies finite incidence structures with prescribed
regularity.

A \(t\)-design is written

\[
\boxed{
t\text{-}(v,k,\lambda).
}
\]

A BIBD is a

\[
\boxed{
2\text{-}(v,k,\lambda)
}
\]

design.

Its parameters satisfy

\[
\boxed{
vr=bk
}
\]

and

\[
\boxed{
r(k-1)
=
\lambda(v-1).
}
\]

Therefore,

\[
\boxed{
r
=
\frac{\lambda(v-1)}{k-1}
}
\]

and

\[
\boxed{
b
=
\frac{\lambda v(v-1)}
{k(k-1)}.
}
\]

Fisher's inequality gives

\[
\boxed{
b\geq v.
}
\]

A symmetric design satisfies

\[
\boxed{
b=v
}
\]

and therefore

\[
\boxed{
r=k.
}
\]

A Steiner system is

\[
\boxed{
S(t,k,v)
=
t\text{-}(v,k,1).
}
\]

Steiner triple systems exist exactly when

\[
\boxed{
v\equiv1
\quad\text{or}\quad
3
\pmod6.
}
\]

A projective plane of order \(q\) is a symmetric design with parameters

\[
\boxed{
2\text{-}(q^2+q+1,q+1,1).
}
\]

An affine plane of order \(q\) is a

\[
\boxed{
2\text{-}(q^2,q,1)
}
\]

design.

Resolvable designs decompose their blocks into parallel classes.

Latin squares and mutually orthogonal Latin squares provide related
regular structures.

Incidence matrices satisfy powerful algebraic identities such as

\[
\boxed{
NN^T
=
(r-\lambda)I+\lambda J.
}
\]

Design theory therefore unifies

\[
\boxed{
\text{counting}
+
\text{finite geometry}
+
\text{algebra}
+
\text{statistics}
+
\text{regular incidence}.
}
\]

%%%%%%%%%%%%%%%%%%%%%%%%%%%%%%%%%%%%%%%%%%%%%%%%%%%%%%%%%%%%
\subsection{Section Connections}
\label{subsec:design-theory-connections}
%%%%%%%%%%%%%%%%%%%%%%%%%%%%%%%%%%%%%%%%%%%%%%%%%%%%%%%%%%%%

\chapterconnection{
Design Theory follows Ramsey Theory by shifting from unavoidable
substructures to deliberately constructed highly regular incidence
structures. Basic Counting supplies the parameter identities, Graph
Theory provides incidence graphs and decomposition interpretations,
Algebraic Combinatorics provides automorphism groups, difference sets,
association schemes, and incidence matrices, while finite-field methods
connect designs with Chapter~\ref{chap:algebra}. Probabilistic and
Extremal Combinatorics contribute packing, covering, and existence
methods. The next section, Matroid Theory, abstracts the concept of
independence from graphs, vector spaces, finite geometries, and designs.
Statistics later connects with the original experimental-design
motivation for balanced block designs.
}

%%%%%%%%%%%%%%%%%%%%%%%%%%%%%%%%%%%%%%%%%%%%%%%%%%%%%%%%%%%%
% SECTION SOURCE TRACKING
%%%%%%%%%%%%%%%%%%%%%%%%%%%%%%%%%%%%%%%%%%%%%%%%%%%%%%%%%%%%

%------------------------------------------------------------
% 6.11 Design Theory
%------------------------------------------------------------
%
% Source basis:
%   - Standard combinatorial design theory.
%   - Standard block-design theory.
%   - Standard finite geometry.
%   - Standard incidence-matrix methods.
%
% Used for:
%   - incidence structures;
%   - t-designs;
%   - BIBDs;
%   - design parameters;
%   - incidence counting;
%   - BIBD parameter identities;
%   - divisibility conditions;
%   - Fisher's inequality;
%   - symmetric designs;
%   - dual designs;
%   - Steiner systems;
%   - Steiner triple systems;
%   - STS existence conditions;
%   - Fano plane;
%   - projective planes;
%   - finite-field projective geometry;
%   - affine planes;
%   - parallel classes;
%   - resolvable designs;
%   - Kirkman's schoolgirl problem;
%   - pairwise balanced designs;
%   - packing designs;
%   - covering designs;
%   - covering numbers;
%   - Latin squares;
%   - orthogonal Latin squares;
%   - mutually orthogonal Latin squares;
%   - experimental block designs;
%   - incidence matrices;
%   - Fisher's inequality by rank;
%   - block intersection graphs;
%   - incidence graphs;
%   - design automorphisms;
%   - difference sets;
%   - group-developed designs;
%   - Hadamard-design preview;
%   - codes from designs;
%   - designs from codes;
%   - orthogonal arrays;
%   - existence problems;
%   - graph decompositions;
%   - hypergraph interpretations;
%   - applications and exercises.
%
% Notes:
%   - Detailed statistical experimental design appears in Chapter 8.
%   - Finite-field constructions connect with Algebra and Number Theory.
%   - Association schemes were previewed in Algebraic Combinatorics.
%   - Matroid Theory follows and generalizes independence structures
%     arising from graphs and finite geometries.
%
%%%%%%%%%%%%%%%%%%%%%%%%%%%%%%%%%%%%%%%%%%%%%%%%%%%%%%%%%%%%